At the correct assumed correlation, DS-RF succeeds in 498, 500 and 499 of 500 asymmetric-mixture attempts at $\rho_{\rm true}=-0.3,0,0.3$, respectively; all Gaussian and OLS-RF attempts succeed. On paired successful mixture samples, the PNIE RMSE ratios (DS-RF/OLS-RF) are 1.011, 0.706 and 0.559, and the interval-length ratios are 1.014, 0.698 and 0.524. DS-RF PNIE coverage is 0.958, 0.952 and 0.946. These results identify how the PNIE precision gain varies with the generating correlation. TNIE RMSE ratios are 0.794, 0.712 and 0.746, with DS-RF coverage between 0.928 and 0.944; TNDE RMSE is higher at the two lower correlations. Total-effect precision and zero-boundary interval lengths are nearly unchanged, and Gaussian results are consistent with the OLS-efficient reference. These patterns connect effect-specific precision gains to curve estimation and distinguish them from estimation of the moment-based zero boundary.

Evaluation at the unadjusted reference $\rho=0$ quantifies the consequences of leaving confounding unaccounted for: for mixture errors at true correlations $-0.3$ and $0.3$, DS-RF PNIE biases are $-0.119$ and $0.120$, and coverage is 0.512 and 0.158. This comparison highlights the complementary roles of nuisance-estimation precision and specification of the sensitivity value in recovering the causal effect. Monte Carlo uncertainty for the 500 attempts is quantified through complete effect-specific standard errors; paired summaries, failures and all-attempted report-and-cover fractions are also supplied in the project repository.
